The proof is in Gauge Theory Gravity with tractors as vectors of $\Cl(2,4)$: adjoin to the spacetime vectors the null vectors $\nO$, $\nInf$ with $\nO \cdot \nInf = 1$, orthogonal to spacetime, and write a tractor field on $U$ as $Y = \sigma\,\nInf + \mu + \rho\,\nO$ with $\mu$ a spacetime vector, so that $\nO \cdot Y = \sigma$, $\nInf \cdot Y = \rho$ and $Q(Y) = Q(\mu) + 2\sigma\rho$, the tractor form of the statement ($\eta^{-1}(\mu, \mu) = Q(\mu)$, a vector $\mu$ standing for the slot $\mu_{a} = a \cdot \mu$). The flat scale is the gravitational field $\hb = 1$, $\Omega = 0$, whose Schouten map vanishes, so the display of Fact \claimref{dark-energy.imported-frozen-point-geometry} reads $\nabla_{a}(\sigma, \mu_{b}, \rho) = (\nabla_{a}\sigma - \mu_{a}, \nabla_{a}\mu_{b} + \eta_{ab}\rho, \nabla_{a}\rho)$; it is the gauge field $D^{T}_{a}Y = a \cdot \nabla Y + \Omega_{T}(a) \times Y$ with $\Omega_{T}(a) = -\nInf \wedge a$, as the computation below shows slot by slot. Since $B \times Y = B \cdot Y$ for a bivector $B$ and $(u \wedge w) \cdot Y = u\,(w \cdot Y) - w\,(u \cdot Y)$,
\[
D^{T}_{a}Y = \bigl(a \cdot \nabla\sigma - a \cdot \mu\bigr)\nInf + \bigl(a \cdot \nabla \mu + \rho\,a\bigr) + \bigl(a \cdot \nabla\rho\bigr)\nO ,
\]
which is the Fact's display in the flat scale, slot by slot.

\emph{(i).} $Y$ is parallel exactly when $a \cdot \nabla\rho = 0$, $a \cdot \mu = a \cdot \nabla\sigma$ and $a \cdot \nabla\mu = -\rho\,a$ for every $a$: on the connected $U$, $\rho = r$ is constant, $\mu = \nabla\sigma$, and $a \cdot \nabla(\nabla\sigma) = -r\,a$. The function $\sigma = s + m \cdot x - \tfrac{r}{2}Q(x)$ solves this ($\nabla\sigma = m - r\,x$), and the difference of two solutions with the same $r$ has constant gradient, so these are all: the parallel tractors are
\[
I_{(s,m,r)}(x) = \sigma(x)\,\nInf + (m - r\,x) + r\,\nO ,
\]
a six-dimensional space. Their norm is $Q(I) = Q(m - r x) + 2r\sigma = Q(m) - 2r\,m \cdot x + r^{2}Q(x) + 2rs + 2r\,m \cdot x - r^{2}Q(x) = 2rs + Q(m)$, constant. The field strength of $\Omega_{T}$ is $(\nInf \wedge a) \times (\nInf \wedge b) = (a \cdot \nInf)\,\nInf \wedge b - (a \cdot b)\,\nInf \wedge \nInf - Q(\nInf)\,a \wedge b + (\nInf \cdot b)\,a \wedge \nInf = 0$ ($\Omega_{T}$ is constant, $\nInf$ null and orthogonal to spacetime): the connection is flat. The Einstein scale is computed directly, not read from the Fact (which states only that an Einstein scale gives a parallel tractor). Where $\sigma \neq 0$ take the gravitational field $\hb^{-1}(a) = a/\sigma$, the dilation of the flat field, whose derived metric is $\sigma^{-2}\eta$. Its torsion equation gives $\Omega(a) = (\nabla\sigma \wedge a)/\sigma$, and with $a \cdot \nabla(\nabla\sigma) = -r\,a$ the field strength is $R(a \wedge b) = -\bigl(2r\sigma + Q(\nabla\sigma)\bigr)\sigma^{-2}\,a \wedge b = -Q(I)\,\sigma^{-2}\,a \wedge b$ (the norm computed above, since $\nabla\sigma = m - r\,x$). As $\hb(B) = \sigma^{2}B$, the store-sign curvature is $\mathcal{R}(B) = -R(\sigma^{2}B) = Q(I)\,B$, so the derived metric $\sigma^{-2}\eta$ is Einstein with $\Lambda = 3Q(I)$. The check computes the same field (rotation-gauge field solved from the torsion condition) and finds $\mathcal{R}(B) = Q(I)\,B = \tfrac{\Lambda}{3}B$.

\emph{(ii).} If $\sigma_{A} = c\,\sigma_{B}$ for a constant $c$ the two quadratic polynomials would be proportional; the check verifies that their coefficient vectors have rank two for the three equal-norm pairs among $I_{A} = (1,0,1)$, $I_{B} = (1,(1,0,0,0),\tfrac12)$, $I_{C} = (2,(0,1,0,0),\tfrac34)$ (at $x = 0$: $\nInf + \nO$, $\nInf + e_{0} + \tfrac12\nO$, $2\nInf + e_{1} + \tfrac34\nO$), so the conformal factor between $\sigma_{A}^{-2}\eta$ and $\sigma_{B}^{-2}\eta$ is not constant. Both derived metrics are Einstein with $\Lambda = 6$ by (i), in the conformal class of $\eta$. The pairing is constant: $\Omega_{T}(a)$ is a bivector, so $a \cdot \nabla(Y \cdot Y') = (D^{T}_{a}Y) \cdot Y' + Y \cdot (D^{T}_{a}Y')$, which vanishes for parallel $Y, Y'$. At $x = 0$, $(\nInf + \nO) \cdot (\nInf + e_{0} + \tfrac12\nO) = \tfrac12 + 1 = \tfrac32$, which is not $\pm 2$, so $I_{B} \neq \pm I_{A}$.

\emph{(iii).} The Lie algebra of $SO(2,4)$ is the space of bivectors of $\Cl(2,4)$, acting by $B \times Y = B \cdot Y$, of dimension $15$. The stabiliser of a non-null $I$ is $\{B : B \cdot I = 0\}$, the bivectors of the five-dimensional $I^{\perp}$, of dimension $10$; that of a pair spanning a nondegenerate plane, $(I_{A} \wedge I_{B})^{2} = (I_{A} \cdot I_{B})^{2} - Q(I_{A})Q(I_{B}) = \tfrac94 - 4 \neq 0$, is the bivectors of the four-dimensional complement of the plane, of dimension $6$. The check computes the three dimensions by exact rank of $B \mapsto (B \cdot I_{A}, B \cdot I_{B})$ on the fifteen basis bivectors. Each frozen point of fixed norm ranges over a five-dimensional quadric, so pairs form a space of dimension $10$; the orbit of the pair has dimension $15 - 6 = 9$, leaving one invariant (the complementary result that where the Weyl part has no kernel vector every parallel tractor is a multiple of one is filed alongside, synthesis.frozen-point-unique-where-weyl-has-no-kernel), and $I_{A} \cdot I_{B}$ is invariant ($B \cdot$ is antisymmetric for the product) and not constant on the space of pairs, so it is that invariant. Proposition \claimref{dark-energy.point-infinity-has-no-physical-fluctuation} concerns one pair $(I, V)$ subject to the Stelle-West constraints, whose fluctuations the nine broken generators exhaust; with a second frozen point the same count leaves the pairing as a physical datum, which that proposition does not address.

\emph{(iv).} For each tower the functional $Q(I_r)^{-1}\int\epsilon(I_r, V_r, F, F)$ is quadratic in the field strength $F$ of the tractor gauge field. Its first variation under $\delta I_r$, $\delta V_r$ (including the variation of the prefactor) is a sum of terms containing $F$ twice, and under a variation $\delta\Omega_{T}$ of the gauge field, with $\delta F = D\,\delta\Omega_{T}$, it is $2Q(I_r)^{-1}\int \epsilon(I_r, V_r, D\,\delta\Omega_{T}, F)$, linear in $F$. Every term vanishes where $F = 0$, which by (i) holds on $U$, so the configuration with any frozen points $I_r$ and any $V_r$ obeying the constraints is stationary for the sum over towers.

\emph{Conclusion.} On $U$ every premise used in (i)-(iv) holds for the two-tower configuration $(I_A, I_B)$, which violates the metric clause of Premise \claimref{dark-matter.p1-single-frozen-point}; so that clause does not follow from them there. Since the parallel tractors on a region of vanishing Weyl part form a six-dimensional space, the pairing takes a continuum of values, and P1 is the further choice $I_B = \pm I_A$, which fixes it at $\pm\Lambda/3$.

\emph{In components} (the flat scale, a reading): $\nabla_a(\sigma, \mu_b, \rho) = (\partial_a\sigma - \mu_a,\ \partial_a\mu_b + \eta_{ab}\rho,\ \partial_a\rho)$, parallel exactly for $\rho = r$, $\mu_a = \partial_a\sigma$, $\partial_a\partial_b\sigma = -r\eta_{ab}$; the norm $\eta^{-1}(d\sigma, d\sigma) + 2\sigma r = 2rs + Q(m)$; and $\mathfrak{so}(2,4)$ as the solutions of $A^{T}G + GA = 0$ for the Gram matrix $G$, which the earlier checks compute.
